% LabHarness demo workspace -- ACS template (achemso class).
%
% Every figure here is generated by a script in scripts/, listed in labharness.toml.
% Run `labharness watch`, change a file in data/, and this document follows.
\documentclass[journal=jacsat,manuscript=article]{achemso}

% LabHarness document style -- generated from the ACS journal style.
%
% One typeface for the whole document, text and figures alike: Latin Modern Roman.
% Do not set fonts anywhere else; change them in the journal style file instead.
\usepackage[T1]{fontenc}
\usepackage{lmodern}
\usepackage{graphicx}
% Editorial tables: rules, numbers aligned on the decimal mark, and notes under the table.
\usepackage{booktabs}
\usepackage{siunitx}
\usepackage{threeparttable}
\sisetup{separate-uncertainty = true}

% achemso loads mciteplus, which initialises its label-width counters at the first
% \bibitem. A document with an empty bibliography never reaches one, and then fails at
% \end{mcitethebibliography} with an undefined control sequence -- which is exactly what
% a freshly created workspace has.
% \gdef, not \providecommand: LaTeX makes \providecommand macros robust, and a robust
% macro does not expand where TeX expects a number.
\makeatletter
\@ifundefined{mcitemaxwidthbibitem}{\gdef\mcitemaxwidthbibitem{0}}{}
\@ifundefined{mcitemaxwidthsubitem}{\gdef\mcitemaxwidthsubitem{0}}{}
\@ifundefined{@mcitecorrectmaxwidthbibitem}{\gdef\@mcitecorrectmaxwidthbibitem{0}}{}
\@ifundefined{@mcitecorrectmaxwidthsubitem}{\gdef\@mcitecorrectmaxwidthsubitem{0}}{}
\makeatother

% \labfigure{path}: include a generated figure at its natural (final) size.
% If the figure does not exist yet, show a placeholder so the document still compiles.
\newcommand{\labfigure}[1]{%
  \IfFileExists{#1}%
    {\includegraphics{#1}}%
    {\fbox{\parbox{0.9\linewidth}{%
      \centering\ttfamily Figure not generated yet:\\\detokenize{#1}%
    }}}%
}

% \labtable{path}: include a generated table. Put it inside a table environment, which
% carries the caption and the label. If it does not exist yet, show a placeholder.
\newcommand{\labtable}[1]{%
  \IfFileExists{#1}%
    {\input{#1}}%
    {\fbox{\ttfamily Table not generated yet: \detokenize{#1}}}%
}

% \labresults{path}: load the macros written next to a fitted plot
% (e.g. figures/calibration.fit.tex) so the text can cite fitted values that update with the data.
\newcommand{\labresults}[1]{\InputIfFileExists{#1}{}{}}

\labresults{figures/calibration.fit.tex}

\title{A demonstration of LabHarness}

\author{First Author}
\email{corresponding.author@example.org}
\affiliation{Department, University, City, Country}

\begin{document}

\begin{abstract}
  A worked example: a catalyst structure, a reaction mechanism and a calibration curve, all
  regenerated from the files in \texttt{data/} whenever they change.
\end{abstract}

\section{Introduction}

Organocatalysis provides a convenient setting for this example.\cite{example2026} The catalyst
is drawn from a single SMILES string, itself produced from a systematic name with
\texttt{labharness resolve}.

\begin{figure}
  \centering
  \labfigure{figures/catalyst.pdf}
  \caption{The catalyst, drawn by RDKit from \texttt{data/catalyst.smi}.}
  \label{fig:catalyst}
\end{figure}

\section{Results and discussion}

The mechanism in Figure~\ref{fig:mechanism} is written in TikZ and chemfig, which is what makes
the curved arrows possible: they are drawn, not approximated.

\begin{figure}
  \centering
  \labfigure{figures/mechanism.pdf}
  \caption{Enamine formation, the first step of the catalytic cycle.}
  \label{fig:mechanism}
\end{figure}

The calibration in Figure~\ref{fig:calibration} was measured in triplicate, so the error bars
are the standard deviation of the replicates rather than a guess. The fitted slope is
\FitCalibrationSlope{} $\pm$ \FitCalibrationSlopeError{} mM$^{-1}$, with
$R^2 = \FitCalibrationRSquared$. That number is not typed here: it is written by the script
that draws the figure, so it follows the data.

\begin{figure}
  \centering
  \labfigure{figures/calibration.pdf}
  \caption{Calibration curve with a linear fit. Error bars are the standard deviation of three
  replicates. The values are demonstration data, not measurements.}
  \label{fig:calibration}
\end{figure}

\section{Conclusions}

Change a number in \texttt{data/calibration.csv} and save: the figure, its error bar and the
slope quoted above all follow, and the PDF reloads on its own.

\bibliography{references}

\end{document}
